Podíl ekonomicky aktivního obyvatelstva pracujícího v~zahraničí vykazuje v~ČR mírně rostoucí trend, přičemž výrazné hodnoty u~SK a~AT dokládají, že přeshraniční pracovní mobilita je v~regionu normou, nikoli výjimkou.
Z~pohledu sociálního dialogu přeshraniční zaměstnávání představuje dvojí výzvu: jednak odliv pracovní síly do vyšší mzdové zóny (AT, DE) vytváří tlak na domácí zaměstnavatele, jednak zaměstnanci pracující v~zahraničí vycházejí z~dosahu domácích \ac{KS} a~odborových organizací.
Rostoucí přeshraniční mobilita je racionální reakcí na mzdový diferenciál (srov.~obr.~\ref{fig:wage_gdp_convergence}): dokud je čistý příjem v~AT nebo DE 1,5--2× vyšší než srovnatelná pozice v~CZ, má odchod ekonomický smysl.
Kolektivní vyjednávání, které by snižovalo tento diferenciál --- prostřednictvím sektorových \ac{KS} ukotvených v~produktivitě --- je tedy nejvýkonnějším nástrojem pro snížení odlivu pracovní síly, a~tím i~pro udržitelnost domácí výrobní kapacity.

Podíl ekonomicky aktivního obyvatelstva pracujícího v~zahraničí vykazuje v~ČR mírně rostoucí trend, přičemž výrazné hodnoty u~SK a~AT dokládají, že přeshraniční pracovní mobilita je v~regionu normou, nikoli výjimkou.
Z~pohledu sociálního dialogu přeshraniční zaměstnávání představuje dvojí výzvu: jednak odliv pracovní síly do vyšší mzdové zóny (AT, DE) vytváří tlak na domácí zaměstnavatele, jednak zaměstnanci pracující v~zahraničí vycházejí z~dosahu domácích \ac{KS} a~odborových organizací.
Rostoucí přeshraniční mobilita je racionální reakcí na mzdový diferenciál (srov.~obr.~\ref{fig:wage_gdp_convergence}): dokud je čistý příjem v~AT nebo DE 1,5--2× vyšší než srovnatelná pozice v~CZ, má odchod ekonomický smysl.
Kolektivní vyjednávání, které by snižovalo tento diferenciál --- prostřednictvím sektorových \ac{KS} ukotvených v~produktivitě --- je tedy nejvýkonnějším nástrojem pro snížení odlivu pracovní síly, a~tím i~pro udržitelnost domácí výrobní kapacity.

Podíl ekonomicky aktivního obyvatelstva pracujícího v~zahraničí vykazuje v~ČR mírně rostoucí trend, přičemž výrazné hodnoty u~SK a~AT dokládají, že přeshraniční pracovní mobilita je v~regionu normou, nikoli výjimkou.
Z~pohledu sociálního dialogu přeshraniční zaměstnávání představuje dvojí výzvu: jednak odliv pracovní síly do vyšší mzdové zóny (AT, DE) vytváří tlak na domácí zaměstnavatele, jednak zaměstnanci pracující v~zahraničí vycházejí z~dosahu domácích \ac{KS} a~odborových organizací.
Rostoucí přeshraniční mobilita je racionální reakcí na mzdový diferenciál (srov.~obr.~\ref{fig:wage_gdp_convergence}): dokud je čistý příjem v~AT nebo DE 1,5--2× vyšší než srovnatelná pozice v~CZ, má odchod ekonomický smysl.
Kolektivní vyjednávání, které by snižovalo tento diferenciál --- prostřednictvím sektorových \ac{KS} ukotvených v~produktivitě --- je tedy nejvýkonnějším nástrojem pro snížení odlivu pracovní síly, a~tím i~pro udržitelnost domácí výrobní kapacity.

Podíl ekonomicky aktivního obyvatelstva pracujícího v~zahraničí vykazuje v~ČR mírně rostoucí trend, přičemž výrazné hodnoty u~SK a~AT dokládají, že přeshraniční pracovní mobilita je v~regionu normou, nikoli výjimkou.
Z~pohledu sociálního dialogu přeshraniční zaměstnávání představuje dvojí výzvu: jednak odliv pracovní síly do vyšší mzdové zóny (AT, DE) vytváří tlak na domácí zaměstnavatele, jednak zaměstnanci pracující v~zahraničí vycházejí z~dosahu domácích \ac{KS} a~odborových organizací.
Rostoucí přeshraniční mobilita je racionální reakcí na mzdový diferenciál (srov.~obr.~\ref{fig:wage_gdp_convergence}): dokud je čistý příjem v~AT nebo DE 1,5--2× vyšší než srovnatelná pozice v~CZ, má odchod ekonomický smysl.
Kolektivní vyjednávání, které by snižovalo tento diferenciál --- prostřednictvím sektorových \ac{KS} ukotvených v~produktivitě --- je tedy nejvýkonnějším nástrojem pro snížení odlivu pracovní síly, a~tím i~pro udržitelnost domácí výrobní kapacity.

\input{texparts/python/cross_border_commuting_timeline}



